% Clase 17 --- El planteo de la integral del solenoide (§2.2).
% "Entonces hagamos un dibujo de costado... voy a tomar el origen en el centro...
%  hago un dibujo más grande para que se vea más claro: Y es la posición de un
%  segmentito, el segmentito tiene ancho dY."
% La figura tiene que dejar clarísimo quién es quién: x es donde se mide, y es
% dónde está la rebanada, y lo que entra en la fórmula de la espira es la
% distancia entre ambos, y - x.
\begin{center}
\begin{tikzpicture}[scale=1]

  % el solenoide de costado
  \draw[figline, line width=1.1pt] (-3.4,-0.95) rectangle (3.4,0.95);
  \draw[figaux] (-4.1,0) -- (4.1,0);
  \node[figlblaux, anchor=west] at (4.15,0) {eje};

  % espiras (de punta): corriente saliente arriba, entrante abajo
  \foreach \xx in {-3.0,-2.4,-1.8,-1.2,-0.6,0,0.6,1.2,1.8,2.4,3.0} {
    \node[figblue, scale=0.62] at (\xx,0.95) {$\odot$};
    \node[figblue, scale=0.62] at (\xx,-0.95) {$\otimes$};
  }

  % origen en el centro
  \node[figgray, circle, fill=figgray, inner sep=1.1pt] at (0,0) {};
  \node[figlblaux, anchor=north] at (0,-0.1) {$0$};

  % la rebanada dy en la posición y
  \fill[figamber!25] (-2.1,-0.95) rectangle (-1.75,0.95);
  \draw[figamber, line width=0.8pt] (-2.1,-0.95) rectangle (-1.75,0.95);
  \draw[figvecaux, figamber] (-2.55,1.55) -- (-2.0,1.05);
  \node[figlbl, figamber, anchor=south east, align=right] at (-2.5,1.5)
    {rebanada de ancho $dy$:\\[-0.25em] contiene $n\,dy$ espiras};
  \node[figlbl, anchor=north] at (-1.93,-2.12) {$y$};

  % el punto de observación
  \node[figred, circle, fill=figred, inner sep=1.5pt] at (1.4,0) {};
  \node[figlbl, figred, anchor=south] at (1.4,0.12) {$P$};
  \node[figlbl, anchor=north] at (1.4,-2.12) {$x$};

  % la distancia que entra en la fórmula
  \draw[figaux] (-1.93,-2.1) -- (-1.93,-0.95);
  \draw[figaux] (1.4,-2.1) -- (1.4,-0.95);
  \draw[figvecaux] (-0.35,-1.5) -- (-1.93,-1.5);
  \draw[figvecaux] (-0.15,-1.5) -- (1.4,-1.5);
  \node[figlbl, anchor=south] at (-0.26,-1.44) {$|y-x|$};

  % cotas L y R
  \draw[figvecaux] (-0.4,1.35) -- (-3.4,1.35);
  \draw[figvecaux] (0.4,1.35) -- (3.4,1.35);
  \node[figlbl, anchor=south] at (0,1.28) {$L$};
  \draw[figvecaux] (3.75,0) -- (3.75,0.95);
  \node[figlbl, anchor=west] at (3.8,0.5) {$R$};

  \node[figlbl, anchor=north, align=center] at (0,-2.8)
    {$dB_z = \dfrac{\mu_0 I R^2}{2\,[(y-x)^2+R^2]^{3/2}}\;n\,dy
     \;\Longrightarrow\;
     B_z = \dfrac{\mu_0 I R^2 n}{2}\displaystyle\int_{-L/2}^{L/2}
     \dfrac{dy}{[(y-x)^2+R^2]^{3/2}}$};

\end{tikzpicture}\\[0.5em]
{\small\itshape El cálculo no agrega física nueva: cada rebanada de ancho $dy$
tiene $n\,dy$ espiras que ---al estar pegadas--- aportan lo mismo que una espira
circular a distancia $y-x$ del punto de observación, y ya sabemos cuánto vale
eso. Conviene fijar el origen en el centro del solenoide y no confundir los dos
lugares: $x$ es dónde se mide y $y$ es dónde está la rebanada que aporta; en la
fórmula de la espira, el ``$Z$'' es la distancia entre los dos. Sustituyendo
$u = y-x$ la integral tiene primitiva elemental,
$u/(R^2\sqrt{u^2+R^2})$, y evaluando entre los extremos queda
$B_z = \frac{\mu_0 n I}{2}\left[\frac{L/2-x}{\sqrt{(L/2-x)^2+R^2}} +
\frac{L/2+x}{\sqrt{(L/2+x)^2+R^2}}\right]$, con los dos términos midiendo
cuánto ``ve'' el punto a cada mitad del solenoide. Ojo con el alcance del
resultado: es exacto (dentro de la aproximación de espiras pegadas) pero
\textbf{sólo sobre el eje}.}
\end{center}
